\section{Categorías cofibradas, categorías bifibradas}
\label{VI.10}
Consideremos una categoría $\cal{F}$ sobre~$\cal{E}$, con el funtor de
proyección
$$
p\colon \cal{F}\to \cal{E},
$$
que define una categoría $\cal{F}^\circ$ sobre $\cal{E}^\circ$ mediante el
funtor de proyección
$$
p^\circ\colon \cal{F}^\circ\to \cal{E}^\circ.
$$
Un morfismo $\alpha\colon \eta\to\xi$ de $\cal{F}$ se llama
\emph{cocartesiano} si es un morfismo cartesiano para $\cal{F}^\circ$
sobre~$\cal{E}^\circ$. Explicitando la definición, se ve que esto significa
que, para todo objeto $\xi'$ de~$\cal{F}_S$, la aplicación
$u\mto u\circ \alpha$
$$
\Hom_S(\xi,\xi')\to\Hom_f(\eta,\xi')
$$
es biyectiva. Se dice también entonces que $(\xi,\alpha)$ es una
\emph{imagen directa}
\index{imagen directa|hyperpage}%
de~$\eta$ por~$f$, en la categoría $\cal{F}$ sobre~$\cal{E}$. Si existe
para todo $\eta$ de~$\cal{F}_T$, se dice que existe el funtor de imagen
directa por~$f$, y se denota este funtor
\marginpar{182}
$f_*^{\cal{F}}$ o~$f_*$,
\label{indnot:ft}\oldindexnot{$f_*^{\cal{F}}$ o~$f_*$|hyperpage}%
una vez escogido. Queda, pues, definido por un isomorfismo de bifuntores
sobre~$\cal{F}_T^\circ\times\cal{F}_S$:
$$
\Hom_S(f_*(\eta),\xi)\isomto \Hom_f(\eta,\xi).
$$
Por consiguiente, si $f_*$ existe, para que $f^*$ exista es necesario y
suficiente que
\ifthenelse{\boolean{orig}}
{$f$}
{$f_*$}
admita un funtor adjunto, es decir, que exista un funtor
$f^*\colon \cal{F}_S\to\cal{F}_T$ y un isomorfismo de bifuntores
$$
\ifthenelse{\boolean{orig}}
{\Hom_E(f_*(\eta),\xi)\isomto \Hom_T(\eta,f^*(\xi)).}
{\Hom_S(f_*(\eta),\xi)\isomto \Hom_T(\eta,f^*(\xi)).}
$$
Sea $g\colon U\to T$ otro morfismo de $\cal{E}$, y supongamos que existen
las imágenes inversas y directas por $f,g$ y~$fg$. Consideremos entonces
las transformaciones naturales
$$
\begin{array}{cccc}
c^{f,g}\colon&f_*g_*&\dpl\from&(fg)_*\\
c_{f,g}\colon&g^*f^*&\dpl\to&(fg)^*.
\end{array}
$$
Se comprueba que, si se consideran $f_*g_*$ y $g^*f^*$ como un par de
funtores adjuntos, al igual que $(fg)_*$ y $(fg)^*$, las dos
transformaciones precedentes son adjuntas entre sí. Por tanto, una es un
isomorfismo si y solo si la otra lo es. En particular:
\begin{proposition}
\label{VI.10.1} Supongamos que la categoría $\cal{F}$ sobre~$\cal{E}$ es
prefibrada y coprefibrada. Para que sea fibrada, es necesario y suficiente
que sea cofibrada.
\end{proposition}

Por supuesto, se dice que $\cal{F}$ es coprefibrada, \resp cofibrada,
\index{cofibrada (categoría)|hyperpage}%
sobre~$\cal{E}$ si $\cal{F}^\circ$ es prefibrada, \resp fibrada%
\ifthenelse{\boolean{orig}}
{, sobre~$\cal{E}$.}
{, sobre~$\cal{E}^\circ$.}
Diremos que $\cal{F}$ es bifibrada sobre~$\cal{E}$,
\index{bifibrada (categoría)|hyperpage}%
si es a la vez fibrada y cofibrada sobre~$\cal{E}$.
